% Printed 62--71, including the intervening unnumbered leaves.
\clearpage
\noindent\makebox[2em][l]{62}\parbox{\dimexpr\linewidth-4em\relax}{\centering\small ANZEIGE.\quad BESTIMMUNG DES BREITENUNTERSCHIEDES ETC.}\hfill\mbox{}\par
\noindent Hr.\ \textsc{Hansen}, nach der mit dem Vf.\ genommenen Abrede gleichzeitig mit diesem an denselben Zenithalsternen mit dem \textsc{Ertel}schen Meridiankreise gemacht hat. Das auf 206 Beobachtungen gegründete Resultat erhält ein noch erhöhtes Interesse durch den Umstand, dass die Seeberger Sternwarte mit dem hannoverschen Dreieckssysteme durch die unter Leitung des Hrn.\ Generallieutenants von \textsc{Müffling} gemessenen Dreiecke in Verbindung ist.

Endlich wird am Schluss noch das Resultat einer auf Veranlassung des Vfs.\ von Hrn.\ Dr.\ \textsc{Schmidt} ausgeführten neuen Berechnung sämmtlicher bisherigen zuverlässigen Gradmessungen mitgetheilt, nach ähnlichen Principien zwar, wie der vorhin erwähnten Arbeit von \textsc{Walbeck} zum Grunde liegen, aber mit schärferer Rechnung, mit Berücksichtigung \textls[150]{sämmtlicher} in den einzelnen Gradmessungen vorkommenden astronomischen Punkte (nicht bloss der Endpunkte), und mit Zuziehung der in gegenwärtigem Werke enthaltenen Resultate der hannoverschen Gradmessung. Man kann diese Bestimmung der Dimensionen des Erdellipsoids als das Zuverlässigste ansehen, was wir bis jetzt aus sämmtlichen Breitengrad-Messungen schliessen können, und es ist interessant, dabei zugleich nach bestimmten Principien unsere Begriffe von dem Grade der Genauigkeit fixirt zu sehen, welcher dadurch bis jetzt erreicht werden kann.
\begin{center}\rule{.25\textwidth}{.4pt}\end{center}

\clearpage
% Unnumbered Bemerkungen opening at the position of printed 63.
\vspace*{14mm}
\begin{center}\textbf{BEMERKUNGEN.}\end{center}
{\small
In dem Originaldrucke des Breitenunterschiedes, der auch durch das noch erhaltene \textsc{Gauss}sche Manuscript controllirt werden konnte, fanden sich in den Zahlenangaben einige kleine Ungenauigkeiten, die nur zum Theil berichtigt wurden.

Geändert wurde in dem vorstehenden Abdrucke auf S.\ 9 die Angabe des Originals für
\begin{center}
\par\noindent\begin{tabular}{ccccccl}
die&Declination&des&Sternes&26:&$50^\circ28'13''{,}12$&in $50^\circ38'13''{,}12$\\
»&»&»&»&40:&$56\quad42\quad21{,}84$&» $56\quad52\quad21{,}74$,
\end{tabular}
\end{center}
da sonst die auf S.\ 22 und 27 mitgetheilten Beobachtungswerthe für diese Sterne mit den nach der Ausgleichung übrig bleibenden Fehlern auf S.\ 35 nicht in Übereinstimmung sind.

In dem II.\ Capitel, den Beobachtungen, sind einigemale kleine Änderungen von 1 bis 2 Hundertstel Secunden in den Endwerthen vorgenommen, um diese mit den Beobachtungswerthen und den Reductionen in Übereinstimmung zu bringen; Voraussetzung war dabei, dass die in die spätere Rechnung eingehenden mittlern Endwerthe hierdurch nicht geändert wurden. Unverbessert mussten die folgenden Stellen bleiben:
\par\noindent\begin{tabular}{@{}p{.065\linewidth}p{.91\linewidth}@{}}
S.\ 21&Bei Stern 25 ($\theta$ Draconis), Altona. Limbus West. Juni 16, folgt aus dem Beobachtungswerthe und den beiden Reductionen der Endwerth um $1''$ grösser, als angegeben ist.\\
S.\ 22&Bei Stern 26, Altona. Limbus Ost. Juni 27, ergibt sich aus dem Beobachtungswerthe und den Reductionen für den Endwerth $\ldots\ 34''{,}98$ anstatt $\ldots\ 34''{,}89$. Der letztere Werth ist jedoch bei der Mittelbildung benutzt worden.\\
S.\ 26&Bei Stern 36, Göttingen. Limbus West, ist das Mittel der Endwerthe $\ldots\ 21''{,}14$, für die Rechnung ist der angegebene Werth $\ldots\ 21''{,}13$ verwendet worden.\\
S.\ 27&Bei Stern 40, Altona. Limbus West. Juni 22, ist der Endwerth, aus dem Beobachtungswerthe und den Reductionen abgeleitet, um $0''{,}10$ grösser, als dort angegeben ist.
\end{tabular}

Beim III.\ Capitel, Resultate, ist geändert worden:
\par\noindent\begin{tabular}{@{}p{.065\linewidth}p{.91\linewidth}@{}}
S.\ 30&der Breitenunterschied aus den Beobachtungen des Sternes 39, für den im Original $2^\circ0'58''{,}51$ steht, in $\ldots\ 58''{,}57$, wie sich aus den Werthen von S.\ 27 ergibt. Der angegebene Werth des Mittels für den Breitenunterschied, S.\ 30, wird dadurch nicht beeinflusst. Das Gewicht des Resultats ist, wie angegeben, $213{,}42$, wofür sich im Original $213{,}41$ findet.\\
S.\ 34&im Art.\ 6, Z.\ 8 v.\ o., die Angabe des Originals für die Verbesserung des Collimationsfehlers in Göttingen: $-0''{,}012$ in $+0''{,}012$, entsprechend den Werthen $3''{,}75$ auf S.\ 32, Z.\ 2 v.\ o., und $3''{,}76$ auf S.\ 34, Z.\ 12 v.\ o. Die Nachrechnung hat die Richtigkeit des $+$ Zeichens bestätigt.
\end{tabular}\par
}

\clearpage
\noindent\makebox[2em][l]{64}\parbox{\dimexpr\linewidth-4em\relax}{\centering\small BEMERKUNGEN.\quad BESTIMMUNG DES BREITENUNTERSCHIEDES ETC.}\hfill\mbox{}\par
{\small
\par\noindent\begin{tabular}{@{}p{.065\linewidth}p{.91\linewidth}@{}}
S.\ 45&in der Tabelle in der Columne für die Fehler
\end{tabular}
\begin{center}
\par\noindent\begin{tabular}{llllll}
der erste&Werth&$-0''{,}31$&des Originals&in&$-0''{,}30$,\\
der letzte&»&$-0{,}23$&»\quad »&»&$-0{,}24$;
\end{tabular}
\end{center}
\begin{quote}
ferner die Summe der Producte der Quadrate der Fehler in die Gewichte, für die im Original zweimal $9{,}6184$ angegeben ist, in $9{,}6154$.
\end{quote}
Nicht geändert ist:
\par\noindent\begin{tabular}{@{}p{.065\linewidth}p{.91\linewidth}@{}}
S.\ 31&bei Stern 9, Altona, das Gewicht des Collimationsfehlers: $4{,}00$, das nach der Anzahl der Beobachtungen auf S.\ 15 gleich $5{,}33$ sein muss, weil mit dem Werthe $4{,}00$ das Endresultat auf S.\ 32 abgeleitet worden ist.
\end{tabular}

Ferner ist zu Art.\ 7, S.\ 35, folgendes zu bemerken. Da sich mit der hier angegebenen Summe der Producte aus den Quadraten der im Art.\ 6 aufgeführten Unterschiede in die entsprechende Anzahl der Beobachtungen: $292{,}8249$ nicht genau der daraus abgeleitete mittlere Fehler einer Beobachtung: $1''{,}5308$ ergab, so wurde eine Nachrechnung dieser Werthe vorgenommen. Diese lieferte an Stelle des ersten Werthes $291{,}8422$ und damit als mittlern Fehler einer Beobachtung $\pm1''{,}5280$ und als mittlern Fehler des Breitenunterschiedes $\pm0''{,}1036$ an Stelle von $\pm0''{,}1038$ des Originals.

Endlich ist im IV.\ Capitel in der Tabelle auf S.\ 53 bei Stern 31 der im Original angegebene Werth der Zenithdistanz: $+4^\circ48'58''{,}28$ in $+4^\circ57'58''{,}28$ umgeändert worden; der letzte Werth ergibt sich rückwärts aus der nebenstehenden Breite und aus der auf S.\ 9 aufgeführten Declination für den Stern.

Die von \textsc{Gauss} in der Anzeige der Bestimmung des Breitenunterschiedes etc., S.\ 61, erwähnte Abhandlung zur lappländischen Gradmessung ist wahrscheinlich die von O.\ A.\ \textsc{Rosenberger}: Über die, auf Veranstaltung der französischen Academie, während der Jahre 1736 und 1737 in Schweden vorgenommene Gradmessung. (Astr.\ Nachr.\ Sechster Band, 1828. S.\ 1--32).
\par\hfill\textsc{Krüger, Börsch.}\par
}
\begin{center}\rule{.25\textwidth}{.4pt}\end{center}

\clearpage
% Unnumbered section title, principal PDF 72 (position 65).
\vspace*{45mm}
\begin{center}
{\Huge ERDELLIPSOID}\par\vspace{10mm}
{\large UND}\par\vspace{10mm}
{\Huge GEODÄTISCHE LINIE.}
\end{center}
\clearpage
% Blank verso, principal PDF 73 (position 66).
\null
\clearpage
% Unnumbered Nachlass opening, principal PDF 74 (position 67).
\begin{center}
{\Large\textls[180]{NACHLASS.}}\par\vspace{12mm}
[1.]\par
Das [Erd]ellipsoid.
\end{center}
Es sei $a$ die halbe grosse, $b$ die halbe kleine Axe; eines Orts astronomische Polhöhe $\varphi$, die sogenannte verbesserte Polhöhe $\varphi'$ [und die reducirte Polhöhe $\psi$]; Abstand vom Mittelpunkt $r$, von der Erdaxe $x$, vom Äquator $y$.

Da sodann vermöge der Gleichung der Ellipse $\dfrac{xx}{aa}+\dfrac{yy}{bb}=1$ ist, so kann man
\[
x=r\cos\varphi'=a\cos\psi,\qquad y=r\sin\varphi'=b\sin\psi
\]
setzen. Nennt man ferner $s$ den elliptischen Bogen zwischen dem Orte und Äquator, so hat man
\[
-\sin\varphi\cdot\mathrm ds=\mathrm dx,\qquad\cos\varphi\cdot\mathrm ds=\mathrm dy,
\]
also
\[
\sin\varphi\cdot\mathrm ds=a\sin\psi\cdot\mathrm d\psi,\qquad\cos\varphi\cdot\mathrm ds=b\cos\psi\cdot\mathrm d\psi.
\]
Hieraus folgt
\[
\operatorname{tang}\psi=\frac ba\operatorname{tang}\varphi,\qquad\operatorname{tang}\varphi'=\frac{bb}{aa}\operatorname{tang}\varphi.
\]
Ferner
\[
\begin{aligned}
\cos\psi&=\frac{a\cos\varphi}{\sqrt{(aa\cos\varphi^2+bb\sin\varphi^2)}}=\sqrt{\frac{\cos\varphi\cos\varphi'}{\cos(\varphi-\varphi')}},\quad&\cos\varphi'&=\frac{aa\cos\varphi}{\sqrt{(a^4\cos\varphi^2+b^4\sin\varphi^2)}}\\
\sin\psi&=\frac{b\sin\varphi}{\sqrt{(aa\cos\varphi^2+bb\sin\varphi^2)}}=\sqrt{\frac{\sin\varphi\sin\varphi'}{\cos(\varphi-\varphi')}},&\sin\varphi'&=\frac{bb\sin\varphi}{\sqrt{(a^4\cos\varphi^2+b^4\sin\varphi^2)}}
\end{aligned}
\]
\[
r=\sqrt{\frac{a^4\cos\varphi^2+b^4\sin\varphi^2}{aa\cos\varphi^2+bb\sin\varphi^2}}=a\sqrt{\frac{\cos\varphi}{\cos\varphi'\cos(\varphi-\varphi')}}=b\sqrt{\frac{\sin\varphi}{\sin\varphi'\cos(\varphi-\varphi')}}
\]
\noindent\hfill $9^*$\par

\clearpage
\noindent 68\hfill NACHLASS.\hfill\mbox{}\par
\[
\frac{a\,\mathrm d\psi}{\cos\psi^2}=\frac{b\,\mathrm d\varphi}{\cos\varphi^2},\quad\text{also}\quad\frac{\mathrm ds}{\mathrm d\varphi}=\frac{bb\cos\psi^3}{a\cos\varphi^3}=\frac{aa\sin\psi^3}{b\sin\varphi^3}=\frac{aabb}{(aa\cos\varphi^2+bb\sin\varphi^2)^{\frac32}}
\]
\[
ee=1-\frac{bb}{aa}.
\]
Ist ein Quadrant des Erdmeridians $=Q$, so ist ein Quadrant des Äquators [vergl.\ Band IV, S.\ 330]
\[
\begin{aligned}
&=Q\cdot(1+\tfrac14ee+\tfrac7{64}e^4+\tfrac{15}{256}e^6+\tfrac{579}{16384}e^8+\tfrac{1515}{65536}e^{10}+\ldots)\\
&=Q:(1-\tfrac14ee-\tfrac3{64}e^4-\tfrac5{256}e^6-\tfrac{175}{16384}e^8-\tfrac{441}{65536}e^{10}-\ldots).
\end{aligned}
\]
\begin{center}\rule{.25\textwidth}{.4pt}\end{center}
Für die Abplattung $\tfrac1{310}$ ist $e=\dfrac{\sqrt{619}}{310}$, $\log e=8{,}9044836$
\[
\begin{array}{l@{\quad}l@{\quad}l}
ee\ldots\ 7{,}808\,9672&e^4\ldots\,\ldotp\ 5{,}617\,9344&e^6\ldots\,\ldotp\,\ldotp\ 3{,}426\,9016\\[4pt]
\tfrac14\ldots\,\ldotp\ 9{,}397\,9400&\tfrac7{64}\ldots\ 9{,}038\,9180&\tfrac{15}{256}\ldots\ 8{,}767\,8513
\end{array}
\]
Quadr.\ des Äqu.\ [wenn $Q=10\,000\,000$ Meter gesetzt wird:]
\[
\begin{array}{r@{\quad}c@{\quad}r}
10\,016\,103^{\mathrm m}{,}015&&\\
45{,}379&&\\
0{,}157&&\\\cline{1-1}
10\,016\,148^{\mathrm m}{,}551&\ldotp\,\ldotp\,\ldotp\,\ldotp\,\ldotp\,\ldotp\,\ldotp&7{,}000\,7008\\
&\left[\dfrac\pi2\right]\ldots&0{,}196\,1199\\\cline{3-3}
&\log a=&6{,}804\,5809.
\end{array}
\]
\begin{center}\rule{.25\textwidth}{.4pt}\end{center}
\[
\text{Abplattung}=\omega,\qquad b=a(1-\omega),\qquad ee=(2-\omega)\omega.
\]
\[
e=\sin E,\qquad\frac{1-\sqrt{(1-ee)}}{1+\sqrt{(1-ee)}}=\operatorname{tang}\tfrac12E^2.
\]
Für $\omega=\dfrac1{302{,}68}$ [ist]
\[
E=4^\circ39'31''{,}355,\qquad\operatorname{tang}\tfrac12E\ldots\ 8{,}609\,3534
\]
\[
\begin{array}{l@{\quad}l@{\quad}l}
ee\ldots\ 7{,}819\,3287&e^4\ldots\,\ldotp\ 5{,}638\,6574&e^6\ldots\,\ldotp\ 3{,}457\,9861\\[4pt]
\tfrac14\ldots\ 9{,}397\,9400&\tfrac7{64}\ldots\ 9{,}038\,9180&\tfrac{15}{256}\ldotp\,\ldotp\ 8{,}767\,8513
\end{array}
\]
[Für $Q=10000000$ Meter ist mithin der Quadrant des Äquators:]

\clearpage
\noindent\mbox{}\hfill DAS ERDELLIPSOID.\hfill 69\par
\[
\begin{array}{r@{\quad}c@{\quad}r}
10\,016\,491^{\mathrm m}{,}8250&&\\
47{,}5954&&\\
0{,}1682&&\\\cline{1-1}
10\,016\,539^{\mathrm m}{,}5886&\ldotp\,\ldotp\,\ldotp\,\ldotp\,\ldotp\,\ldotp\,\ldotp\,\ldotp\,\ldotp\,\ldotp\,\ldotp\,\ldotp\,\ldotp\,\ldotp&7{,}000\,7177\,118\\
&\left[\dfrac\pi2\right]\ldots&0{,}196\,1198\,770\\\cline{3-3}
\multicolumn{2}{r}{\log a\text{ (in Metern)}=}&6{,}804\,5978\,348
\end{array}
\]
[Da der] Meridianquadrant $=5130878{,}3$ Toisen [ist, so ist]
\[
\log a\text{ (in Toisen)}=6{,}514\,7893\,106.
\]
\begin{center}\rule{.25\textwidth}{.4pt}\end{center}
Zahlen, das Erdsphäroid betreffend.
\[
\text{Abplattung}=\frac1{302{,}78},\quad ee=\frac{604{,}56}{(302{,}78)^2},\quad\log ee=7{,}81918\,50399{\scriptstyle45}.
\]
[Der Umfang des Meridians ist $=$] Peripherie des Äquators mal
\[
\begin{aligned}
\bigg(1-\tfrac14ee&-\frac{1.3}{4.16}e^4-\frac{1.3.15}{4.16.36}e^6-\frac{1.3.15.35}{4.16.36.64}e^8-\frac{1.3.15.35.63}{4.16.36.64.100}e^{10}\\
&\bigg[-\frac{1.3.15.35.63.99}{4.16.36.64.100.144}e^{12}-\frac{1.3.15.35.63.99.143}{4.16.36.64.100.144.196}e^{14}\bigg]-\ldots\bigg).
\end{aligned}
\]
\[
\begin{array}{rcl}
\tfrac14ee&=&\multicolumn{1}{r}{0{,}00164\,86370\,23537\,4099}\\[3pt]
\tfrac3{64}e^4&=&\multicolumn{1}{r}{20385\,03026\,5312}\\[3pt]
\tfrac5{256}e^6&=&\multicolumn{1}{r}{56\,01252\,6094}\\[3pt]
\tfrac{175}{16384}e^8&=&\multicolumn{1}{r}{20200\,3209}\\[3pt]
\tfrac{441}{65536}e^{10}&=&\multicolumn{1}{r}{83\,9236}\\[3pt]
\tfrac{4851}{1048576}e^{12}&=&\multicolumn{1}{r}{3805}\\[3pt]
\left[\tfrac{693693}{205520896}e^{14}\right]&=&\multicolumn{1}{r}{18}\\\cline{3-3}
&&0{,}00165\,06811\,48101\,1773\\[3pt]
\text{Summe der Reihe}&=&0{,}99834\,93188\,51898\,8227
\end{array}
\]
\[
\begin{array}{rcl}
\text{Logarithm.}&\ldots&9{,}99928\,25259{\scriptstyle92}\\
\dfrac{20\,000\,000}{\pi}&\ldotp\,\ldotp\,\ldotp\,\ldotp\,\ldotp\,\ldotp&6{,}80388\,01229{\scriptstyle70}\\\cline{3-3}
\text{\raisebox{.5ex}{.}}\quad\log a&=&6{,}80459\,75969{\scriptstyle78}.
\end{array}
\]

\clearpage
\noindent 70\hfill NACHLASS.\hfill\mbox{}\par
\begin{center}[2.]\par
[Gleichung der Verticalebene des Rotationsellipsoids.]\end{center}
Für die Örter im Schnitt der Oberfläche des Ellipsoids mit einer Ebene, die bei der Länge $0$, Polhöhe $\varphi^0$ auf jener senkrecht ist und den nördlichen Theil des dortigen Meridians unter dem Winkel $\zeta^0$ schneidet, ist die Bedingungsgleichung diese:
\[
\frac{\cos\varphi\cos\lambda\sin\varphi^0}{\sqrt{(1-ee\sin\varphi^2)}}+\frac{\cos\varphi\sin\lambda\operatorname{cotang}\zeta^0}{\sqrt{(1-ee\sin\varphi^2)}}-\frac{(1-ee)\sin\varphi\cos\varphi^0}{\sqrt{(1-ee\sin\varphi^2)}}=\frac{ee\cos\varphi^0\sin\varphi^0}{\sqrt{(1-ee\sin\varphi^{0^2})}}
\]
oder einfacher
\[
\cos\varphi\cos\lambda\sin\varphi^0+\cos\varphi\sin\lambda\operatorname{cotang}\zeta^0-(1-ee)\sin\varphi\cos\varphi^0=ee\cos\varphi^0\sin\varphi^0\sqrt{\frac{1-ee\sin\varphi^2}{1-ee\sin\varphi^{0^2}}}.
\]
\begin{center}[3.]\end{center}
Gleichung des [Rotations]ellipsoids in Beziehung auf eine berührende Ebene:
\[
xx(1-ee\sin\varphi^2)+yy(1-ee)+xzee\sin2\varphi+zz(1-ee\cos\varphi^2)-\frac{2bbz}{a\sqrt{(1-ee\sin\varphi^2)}}=0.
\]
[Die $x$-Axe ist Tangente der Meridianellipse, positiv nach Süden; die $z$-Axe fällt mit der Richtung der Normalen zusammen und ist positiv nach dem Innern des Ellipsoids.]

[Man hat auch:]
\[
\left(\frac{x\sin\varphi-z\cos\varphi}{a}\right)^2+\left(\frac{x\cos\varphi+z\sin\varphi}{b}\right)^2+\frac{yy}{aa}=\frac{2z}{a\sqrt{(1-ee\sin\varphi^2)}}.
\]
[Setzt man]
\[
\begin{aligned}
x&=s\cos\delta\cos\theta\\
y&=s\cos\delta\sin\theta\\
z&=s\sin\delta,
\end{aligned}
\]
[wo $\delta$ der Depressionswinkel und $\theta$ das Azimuth ist, so wird]
\[
s\left(1+\frac{ee}{1-ee}(\cos\varphi\cos\theta\cos\delta+\sin\varphi\sin\delta)\right)^2\sqrt{(1-ee\sin\varphi^2)}=2a\sin\delta.
\]
[Angenähert ist]
\[
\sin\lambda=\frac{s\cos\delta\sin\theta\sqrt{(1-ee\sin\varphi^2)}}{a\cos\varphi}.
\]
\begin{center}\rule{.18\textwidth}{.4pt}\end{center}

\clearpage
\noindent\mbox{}\hfill DAS ERDELLIPSOID.\hfill 71\par
\vspace{12mm}
\begin{center}\textbf{BEMERKUNGEN.}\end{center}
{\small
Die Notiz [1] ist einem Handbuche entnommen; die dazu gehörigen Zahlenwerthe fanden sich in zwei andern Handbüchern; die Zahlen, das Ellipsoid betreffend, sind auf die letzte Seite des \textsc{Gauss}schen Exemplars der »Mathematischen Abhandlungen von Dr.\ H.\ F.\ \textsc{Scherk}. Berlin, 1825« eingetragen. Der Abplattungswerth $\tfrac1{310}$ war der vom General von \textsc{Müffling} bei den preussischen Vermessungen zu Grunde gelegte. Wie aus dem Briefe an \textsc{Olbers} vom 18. April 1822 hervorgeht, hat \textsc{Gauss} im Anfange seiner Gradmessungsarbeiten aus Versehen an Stelle des \textsc{Walbeck}schen Werthes die Abplattung $\dfrac1{302{,}68}$ benutzt. In der That finden sich verschiedene Tabellen und Rechnungen, die hierauf beruhen. \textsc{Gauss}, der anfangs diesen Werth beibehalten wollte, entschloss sich jedoch später (vergleiche den Brief von \textsc{Gauss} an \textsc{Olbers} vom 1. März 1827) zu \textsc{Walbeck}s Werth $\dfrac1{302{,}78}$ überzugehen. Dieser liegt also der Berechnung der \textsc{Gauss}schen Messungen zu Grunde. Die Länge des Meridianquadranten ist nach \textsc{Walbeck} $5130878{,}22$ Toisen (Dissertatio de forma et magnitudine telluris ex dimensis arcubus meridiani definiendis. Aboae, 1819).

[2] und der erste Theil von [3] gehören demselben Handbuche wie [1] an; der zweite Theil von [3] fand sich auf einem einzelnen Blatte.
\par\hfill\textsc{Krüger, Börsch.}\par
}
\begin{center}\rule{.25\textwidth}{.4pt}\end{center}
